\title{Controlling Text-to-Image Diffusion by Orthogonal Finetuning}

\begin{abstract}
Large-scale text-to-image diffusion models have demonstrated remarkable abilities in synthesizing photorealistic visuals based on textual descriptions. However, developing effective strategies to direct or tailor these advanced models for specific downstream applications remains a significant and unresolved issue. To address this, we present a structured finetuning approach—Orthogonal Finetuning~(OFT)—to adapt text-to-image diffusion models for diverse downstream tasks. In contrast to prior approaches, OFT is theoretically guaranteed to maintain hyperspherical energy, which encodes the inter-neuronal relationships on a unit hypersphere. Our analysis shows that retaining this characteristic is essential for sustaining the semantic generative potential of text-to-image diffusion models. For further stability during the finetuning process, we propose Constrained Orthogonal Finetuning (COFT), which incorporates an additional constraint on the radius of the hypersphere. In particular, our study examines two core finetuning scenarios in text-to-image generation: subject-driven generation—where the aim is to synthesize images of a specific subject, given a small set of subject images and a textual description; and controllable generation—where the objective is to generate outputs conditioned on extra control cues. Experimental evidence indicates that the OFT framework surpasses contemporary alternatives both in terms of image quality and finetuning speed.
\end{abstract}

\section{Introduction}

Advances in text-to-image diffusion models~\cite{saharia2022photorealistic,ramesh2022hierarchical,rombach2022high} have resulted in strong capabilities for text-conditioned, high-quality image synthesis. Yet, the supervision provided by text prompts may remain vague or fall short in delivering the precise, detailed control desired over generated images. In response, this work focuses on two principal tasks in the text-to-image generation domain:

\begin{itemize}[leftmargin=*,nosep]

    \item \textbf{Subject-driven generation}~\cite{ruiz2023dreambooth}: Starting with a limited number of images featuring a particular subject, the challenge is to synthesize further images of that subject situated in new environments, guided by textual input.
    \item \textbf{Controllable generation}~\cite{zhang2023adding,mou2023t2i}: When presented with supplementary control information (\emph{e.g.}, canny edge maps or segmentation cues), the objective is to generate images conforming both to this control input and to a textual prompt.
\end{itemize}

The core challenge in both tasks lies in how to efficiently finetune text-to-image diffusion models while maintaining the generative strengths attained during pretraining. An ideal finetuning approach should fulfill two primary requirements: (1) \emph{training efficiency}, involving minimal trainable parameters and a reduced number of training epochs, and (2) \emph{generalizability preservation}, which ensures the model retains its capacity for producing high-fidelity, varied outputs. Common strategies for finetuning involve either making small adjustments to the existing neuron weights using a modest learning rate (\emph{e.g.}, \cite{ruiz2023dreambooth}) or incorporating lightweight modules that re-parameterize certain weights (\emph{e.g.}, \cite{hulora2022,zhang2023adding}). However, despite their practicality, these approaches do not inherently guarantee the conservation of generative abilities learned during pretraining. Furthermore, there is a limited foundational understanding guiding the development of robust finetuning procedures and the selection of key hyperparameters, such as the appropriate number of training epochs. A significant obstacle is the absence of a concrete metric for evaluating how well the generative properties of the pretrained model are preserved post-finetuning. Current approaches generally assume that a smaller Euclidean distance between the pretrained and finetuned models corresponds to better retention of pretrained capabilities, which motivates the adoption of very low learning rates. Nonetheless, we contend that this reliance on Euclidean proximity is insufficient for thoroughly assessing semantic preservation. Introducing a more structurally informative metric to characterize changes between the finetuned and original models could substantially improve both the stability of the finetuning process and the retention of pretraining performance.

Building upon previous findings that hyperspherical similarity effectively captures semantic relationships~\cite{liu2017deep,liu2018decoupled,chen2020angular}, we utilize hyperspherical energy~\cite{liu2018learning} as a means to represent the pairwise interactions among neurons. Hyperspherical energy, calculated as the aggregate of hyperspherical similarities (such as cosine similarity) for every pair of neurons in a particular layer, measures how uniformly neurons are distributed across the unit hypersphere~\cite{liu2021learning}. We put forward the conjecture that an effective finetuned model should exhibit only minor changes in hyperspherical energy relative to its pretrained counterpart. A straightforward solution might be to employ a regularization term to stabilize hyperspherical energy during finetuning; however, this approach does not necessarily ensure that the energy difference will be minimal. To address this, we exploit a key invariance of hyperspherical energy: when the same orthogonal transformation is applied to all neurons, their pairwise hyperspherical similarities are mathematically invariant. Guided by this property, we introduce Orthogonal Finetuning~(OFT), a technique for adapting large text-to-image diffusion models to downstream tasks while preserving their hyperspherical energy. The foundational idea behind OFT is to train a single orthogonal transformation per layer to maintain the angular relationships among neurons. Conceptually, OFT can also be understood as redefining the canonical basis for neurons within each layer. Further considering that a lower Euclidean distance between the finetuned and pretrained models is correlated with retention of the original performance, we advance an enhanced version—Constrained Orthogonal Finetuning~(COFT)—which restricts the finetuned neurons to remain within a hypersphere of fixed radius around the original pretrained neurons.

The rationale behind employing orthogonal transformation for neuron finetuning is partly informed by insights from the 2D Fourier transform, wherein an image can be represented through its magnitude and phase spectra. Notably, the phase spectrum—reflecting the angular relationship between inputs and basis functions—retains the principal semantic content. For instance, combining the phase spectrum from an image with a randomly chosen magnitude spectrum typically suffices to reconstruct the original image while preserving its semantics. This observation implies that adjusting neuron orientations predominantly influences the semantic characteristics of generated images, an objective at the core of both subject-driven and controllable generation. Nonetheless, unconstrained alteration of neuron directions risks compromising the pretrained model’s generative capabilities. To address this, we advocate for maintaining the inter-neuronal angles, predicated on the assumption that these angles are vital for encoding neural network knowledge. Building on this perspective, we introduce a shared orthogonal transformation across layers to ensure that the hyperspherical energy remains invariant.

Additionally, our approach draws from orthogonal over-parameterized training~\cite{liu2021orthogonal}, a method which trains classification networks by applying orthogonal transformations to randomly initialized models. Since random initialization typically results in a neural network with low expected hyperspherical energy, \cite{liu2021orthogonal} focuses on preserving this low energy throughout training—an objective linked to improved generalization in classification tasks~\cite{liu2018learning,lin2020regularizing}. Their work demonstrates that orthogonal transformation grants sufficient capacity for building well-generalizing classification models. Conversely, our primary focus lies in finetuning text-to-image diffusion models to enhance controllability and maintain, or improve, downstream generative performance. There are two primary distinctions between our Orthogonal Finetuning Transformation (OFT) and the strategy in \cite{liu2021orthogonal}. First, unlike \cite{liu2021orthogonal}, which aims to further minimize hyperspherical energy, OFT seeks to conserve the hyperspherical energy inherent to the pretrained model, thereby preserving its essential structure during finetuning; indiscriminate energy minimization in diffusion models risks degrading their underlying semantic architecture. Second, OFT introduces explicit constraints to minimize divergence from the pretrained state, a consideration absent in \cite{liu2021orthogonal}. Effective finetuning thus hinges on carefully balancing adaptability with preservation of learned features, a trade-off our OFT method is specifically designed to optimize. Below, we summarize our main contributions:

\begin{itemize}[leftmargin=*,nosep]

    \item We introduce Orthogonal Finetuning (OFT), a new approach for adapting text-to-image diffusion models to enhance their controllability. To address concerns about training instability, we further devise a constrained form of OFT that restricts the model's angular deviation from its original pretrained state.
    \item Unlike existing finetuning techniques, OFT updates model parameters while analytically guaranteeing the conservation of hyperspherical energy—a property we empirically identify as critical for retaining the pretrained model's generative semantic capacity.
    \item We evaluate OFT on both subject-driven and controllable image generation tasks. Through a thorough set of experiments, we establish that OFT consistently surpasses previous methods in terms of generated image quality, speed of convergence, and training robustness. Furthermore, OFT demonstrates improved data efficiency, effectively achieving convergence with substantially fewer training images and epochs.
    \item For the controllable generation setting, we propose a novel control consistency metric to quantitatively measure controllability. The key concept is to infer the control input from the model output and compare it to the original control specification. This metric further confirms OFT's high level of controllability.
\end{itemize}

\section{Related Work}

\textbf{Text-to-image diffusion models}. Rapid advances have recently been achieved in text-to-image synthesis~\cite{saharia2022photorealistic,ramesh2022hierarchical,rombach2022high,gu2022vector,nichol2022glide}, fueled by innovations in diffusion-based generative modeling~\cite{ho2020denoising,song2021score,song2021denoising,dhariwal2021diffusion} as well as vision-language joint representations~\cite{su2020vlbert,lu2019vilbert,radford2021learning,wang2021simvlm,li2022blip,li2023blip,alayrac2022flamingo,schuhmann2022laion}. Notably, models like GLIDE~\cite{nichol2022glide} and Imagen~\cite{saharia2022photorealistic} learn in pixel space: GLIDE jointly optimizes the text encoder with a diffusion model using aligned text-image pairs, while Imagen relies on a fixed pretrained text encoder throughout training. Conversely, Stable Diffusion~\cite{rombach2022high} and DALL-E2~\cite{ramesh2022hierarchical} employ latent space diffusion models. Stable Diffusion utilizes VQ-GAN~\cite{esser2021taming} to obtain a latent visual codebook, whereas DALL-E2 incorporates CLIP~\cite{radford2021learning} to generate a shared latent embedding space for multimodal representation. Beyond diffusion-based approaches, the community has also explored generative adversarial networks~\cite{reed2016generative,zhang2017stackgan,xu2018attngan,li2019controllable} and autoregressive strategies~\cite{ramesh2021zero,ding2021cogview,wu2022nuwa,yuscaling} for this task. OFT is designed to be independent of model architecture and thus can be applied to any diffusion-based text-to-image model.

\textbf{Subject-driven generation}. To avoid undesired changes to the target subject, works such as~\cite{avrahami2022blended,nichol2022glide} incorporate user-provided masks as supplementary constraints. Inversion-based techniques~\cite{choi2021ilvr,dhariwal2021diffusion,ramesh2022hierarchical,gal2022image} facilitate context manipulation while retaining subject features, and methods like~\cite{hertz2022prompt} support editing both local and global image regions without relying on explicit masks. However, these strategies may struggle to preserve fine identity details. Approaches such as Pivotal Tuning~\cite{roich2022pivotal} adapt generators by finetuning around an inverted latent point with regularization to maintain subject identity, while~\cite{nitzan2022mystyle} trains a person-specific generative face model using a personal image dataset. In~\cite{casanova2021instance}, instance-based variations are generated, yet may fail to fully capture instance-specific attributes. DreamBooth~\cite{ruiz2023dreambooth} leverages custom textual tokens and a handful of subject images to update text-to-image diffusion models using both reconstruction and class-specific prior preservation losses. Building upon the DreamBooth procedure, OFT applies orthogonal updates in place of conventional finetuning with small learning rates, aiming to preserve pretrained semantics while improving adaptation.

\textbf{Controllable generation}. Image-to-image translation is typically regarded as a controllable generation problem, where most prior techniques leverage conditional generative adversarial networks~\cite{isola2017image,zhu2017toward,wang2018high,park2019semantic,choi2018stargan}. In addition, diffusion models have also been applied to this domain~\cite{saharia2022palette,voynov2022sketch,wang2022pretraining}. More recently, the introduction of ControlNet~\cite{zhang2023adding} enables enhanced control over pretrained diffusion models by fine-tuning them with supplementary control signals, resulting in state-of-the-art controllable image synthesis. Similarly, T2I-Adapter~\cite{mou2023t2i} also achieves increased control by fine-tuning a pretrained diffusion model, aligning with the approach adopted in concurrent research. In line with these frameworks~\cite{zhang2023adding,mou2023t2i}, our work utilizes OFT to fine-tune pretrained diffusion models, consistently producing superior controllability even with reduced training samples and fewer trainable parameters. Notably, OFT also ensures that no additional computational burden is introduced during inference.

\textbf{Model finetuning}. The strategy of fine-tuning large pretrained models on specific downstream applications has gained significant momentum in recent years~\cite{devlin2018bert,brown2020language,he2022masked}. Various adaptation techniques, as forms of fine-tuning (\emph{e.g.}, \cite{hulora2022,houlsby2019parameter,pfeiffer2020adapterfusion}), have been extensively explored in the field of natural language processing. The method most directly related to OFT is LoRA~\cite{hulora2022}, which imposes a low-rank constraint on additive weight modifications during fine-tuning. In contrast, OFT applies a layer-consistent orthogonal transformation to neuron weights, updating them multiplicatively. This approach guarantees preservation of the angular relationships between neurons within a layer, offering enhanced training stability.

\section{Orthogonal Finetuning}

\subsection{Why Does Orthogonal Transformation Make Sense?}

To motivate the use of orthogonal transformations in the context of fine-tuning text-to-image diffusion models, we break down the rationale into two principal considerations: (1) the importance of modifying neuron angles (i.e., their directions) during fine-tuning, and (2) the suitability of orthogonal transformations for implementing such angle modifications.

To address the first question, we take cues from empirical findings reported in \cite{liu2018decoupled,chen2020angular}, which highlight that differences in angular features serve as an effective indicator of the semantic gap. SphereNet~\cite{liu2017deep} demonstrates that enforcing normalization of all neural activations onto a unit hypersphere during network training can preserve expressive power and enhance generalization capabilities. This suggests that neuron orientation alone is sufficient to encode the most salient information in the data. To further illustrate the significance of neuron orientation, we present a toy experiment depicted in Figure~\ref{fig:toy_ae}, where we train a conventional convolutional autoencoder on a set of flower images. During training, the model employs a standard inner product for feature map generation ($z$ refers to an output element derived from the convolution kernel $\bm{w}$ applied to input $\bm{x}$ within a moving window). For evaluation, we compare three approaches to compute the feature map: (a) repeating the inner product as used during training, (b) isolating magnitude information, and (c) retaining only angular information. As shown in Figure~\ref{fig:toy_ae}, angular information alone nearly reconstructs the original images, whereas using only magnitude results in a loss of meaningful information. Notably, the network is never trained with cosine activation (c); all training is performed with inner products. These observations indicate that neuron orientation (direction) encodes most of the semantic content present in the input images. Consequently, targeting neuron directions for finetuning should be particularly effective for semantic modifications.

Regarding the second question, the most straightforward approach to adjust neuron directions is to apply a unified rotation or reflection to all neurons in a given layer, inherently involving an orthogonal transformation. While it is possible to devise more adaptable angular transformations—such as rotating individual neurons by distinct angles—we find that orthogonal transformations strike a valuable balance between adaptability and structural regularity. Further, \cite{liu2021orthogonal} provides evidence that orthogonal transformations possess sufficient expressivity for neural network learning tasks. To empirically reinforce this claim, we present an experiment designed to demonstrate the regularizing influence of orthogonality. Employing the controllable generation framework of ControlNet~\cite{zhang2023adding}, we present the outcomes in Figure~\ref{fig:ortho}. The findings show that our standard OFT approach maintains stable performance and delivers precise control after convergence (at epoch 20), whereas an OFT variant lacking orthogonality fails to yield realistic images or demonstrate control. These results underscore the necessity of enforcing orthogonality in OFT.

\subsection{General Framework}

Standard approaches to finetuning often rely on gradient descent with a relatively low learning rate to adjust the model, or selected layers, from its pretrained state. By employing a small learning rate, these methods implicitly restrict how much the finetuned model deviates from the original, reflecting an underlying goal of standard finetuning: to optimize model parameters while implicitly minimizing $\thickmuskip=2mu \medmuskip=2mu \| \bm{M} - \bm{M}^0\|$, where $\bm{M}$ and $\bm{M}^0$ denote the finetuned and pretrained weights, respectively. Although this implicit regularization is intuitive, it may not be stringent enough, especially when finetuning large-scale models. To enhance this, LoRA imposes a further low-rank constraint on weight modifications, specifically $\thickmuskip=2mu \medmuskip=2mu \text{rank}(\bm{M} - \bm{M}^0)=r'$, where $r'$ is set to a low value. In contrast, OFT takes a different approach by constraining the pairwise neuron relationships through hyperspherical energy, such that $\thickmuskip=2mu \medmuskip=2mu \| \text{HE}(\bm{M})-\text{HE}(\bm{M}^0) \|=0$, where $\text{HE}(\cdot)$ denotes the hyperspherical energy for a matrix. 

To clarify, consider a fully connected layer given by $\thickmuskip=2mu \medmuskip=2mu \bm{W}=\{\bm{w}_1,\cdots,\bm{w}_n\}\in\mathbb{R}^{d\times n}$, where each $\bm{w}_i\in\mathbb{R}^{d}$ represents an individual neuron, and $\bm{W}^0$ refers to the pretrained parameters. The output vector from this layer, $\thickmuskip=2mu \medmuskip=2mu \bm{z}\in\mathbb{R}^n$, is calculated as $\thickmuskip=2mu \medmuskip=2mu \bm{z}=\bm{W}^\top\bm{x}$ for input $\thickmuskip=2mu \medmuskip=2mu \bm{x}\in\mathbb{R}^d$. The objective of OFT is then to minimize discrepancies in hyperspherical energy between finetuned and pretrained versions:

\begin{equation}\label{eq:oft_principle}
\footnotesize
\min_{\bm{W}} \left\| \text{HE}(\bm{W})-\text{HE}(\bm{W}^0)\right\|~~\Leftrightarrow~~\min_{\bm{W}} \bigg{\|} \sum_{i\neq j} \|\hat{\bm{w}}_i-\hat{\bm{w}}_j\|^{-1}-\sum_{i\neq j} \|\hat{\bm{w}}^0_i-\hat{\bm{w}}^0_j\|^{-1}\bigg{\|}
\end{equation}

Here, $\thickmuskip=2mu \medmuskip=2mu \hat{\bm{w}}_i:=\bm{w}_i/\|\bm{w}_i\|$ represents the normalization of neuron $i$, and the hyperspherical energy for a layer $\bm{W}$ is given by $\thickmuskip=2mu \medmuskip=2mu \text{HE}(\bm{W}):= \sum_{i\neq j} \|\hat{\bm{w}}_i-\hat{\bm{w}}_j\|^{-1}$. It is apparent from Eq.~\eqref{eq:oft_principle} that the minimum possible value is zero. This zero is attainable whenever $\bm{W}$ is derived from $\bm{W}^0$ via a rotation or reflection; formally, $\thickmuskip=2mu \medmuskip=2mu \bm{W}=\bm{R}\bm{W}^0$ for some orthogonal matrix $\thickmuskip=2mu \medmuskip=2mu \bm{R}\in\mathbb{R}^{d\times d}$, where the determinant of $\bm{R}$ being $1$ or $-1$ signifies a rotation or reflection, respectively. The essence of OFT is to confine finetuning to

\begin{equation}\label{eq:oft_general}
    \footnotesize
    \bm{z}=\bm{W}^\top\bm{x}=(\bm{R}\cdot\bm{W}^0)^\top\bm{x},~~~\text{s.t.}~\bm{R}^\top\bm{R}=\bm{R}\bm{R}^\top=\bm{I}
\end{equation}

Here, $\bm{W}$ represents the weight matrix adapted via OFT, while $\bm{I}$ stands for the identity matrix. An overview of the OFT approach is provided in Figure~\ref{fig:block}. To parallel the zero initialization strategy used in LoRA, it is necessary for OFT to commence fine-tuning from the pretrained weights $\bm{W}^0$. This is ensured by setting the initial value of the orthogonal matrix $\bm{R}$ to the identity matrix, thereby aligning the starting point of the fine-tuned and pretrained models. In order to preserve the orthogonality of $\bm{R}$ during training, differentiable orthogonalization techniques, such as those proposed by \cite{liu2021orthogonal,lezcano2019cheap}, may be employed. Later in the text, we will elaborate on efficient methods for maintaining matrix orthogonality.

\subsection{Efficient Orthogonal Parameterization}\label{sect:eff_ortho}

Although conventional methods like the differentiable Gram-Schmidt process can enforce orthogonality, they are often computationally impractical for large-scale models~\cite{liu2021orthogonal}. To address scalability concerns, we utilize the Cayley transform for constructing orthogonal matrices. Specifically, this involves defining $\thickmuskip=2mu \medmuskip=2mu \bm{R}=(\bm{I}+\bm{Q})(I-\bm{Q})^{-1}$, where $\bm{Q}$ is a skew-symmetric matrix such that $\thickmuskip=2mu \medmuskip=2mu \bm{Q}=-\bm{Q}^\top$. While the Cayley transform offers improved efficiency, it is limited to generating orthogonal matrices with determinant equal to $1$, thereby restricting $\bm{R}$ to the special orthogonal group. However, this constraint has not demonstrated adverse effects on model performance in practice. Even with the Cayley parameterization, using a full orthogonal matrix $\bm{R}$ can remain parameter-intensive when the dimension $d$ is large. To reduce parameter overhead, we suggest structuring $\bm{R}$ as a block-diagonal matrix comprised of $r$ orthogonal submatrices, yielding:

\begin{equation}
    \footnotesize
    \bm{R}=\text{diag}(\bm{R}_1,\bm{R}_2,\cdots,\bm{R}_r)=\begin{bmatrix}
    \bm{R}_1\in \text{O}(\frac{d}{r}) &  &\\
    &\ddots & \\
    & & \bm{R}_r\in \text{O}(\frac{d}{r})
    \end{bmatrix}\in \text{O}(d)
\end{equation}

Here, $\text{O}(d)$ represents the group of $d$-dimensional orthogonal matrices. Both $\bm{R}\in\mathbb{R}^{d\times d}$ and each $\bm{R}_i\in\mathbb{R}^{d/r\times d/r}$, for all $i$, are defined as orthogonal matrices. When $r=1$, the block-diagonal orthogonal matrix reduces to a regular, unconstrained orthogonal matrix. For an orthogonal matrix with dimensions $d\times d$, there are $d(d-1)/2$ independent parameters, giving a parameter complexity of $\mathcal{O}(d^2)$. By contrast, for a block-diagonal orthogonal matrix with $r$ blocks, the parameter count reduces to $d(d/r-1)/2$, i.e., complexity $\mathcal{O}(d^2/r)$. An additional reduction is possible by tying the blocks together—specifically, setting $\bm{R}_i=\bm{R}_j$ for all $i\neq j$—leading to parameter complexity $\mathcal{O}(d^2/r^2)$. Despite these parameter-reduction strategies, the resulting $\bm{R}$ still maintains orthogonality, thereby ensuring that the hyperspherical energy is unchanged.

To analyze parameter efficiency, we contrast OFT with LoRA. For LoRA with low-rank $r'$, the number of learnable parameters is $r'(d+n)$. If $r$ and $r'$ are both defined as a fraction of the neuron dimension (e.g., $r = r' = \alpha d$ for some constant $0<\alpha\leq 1$), then LoRA’s parameter complexity is $\mathcal{O}(d^2 + dn)$, whereas for OFT, it reduces to $\mathcal{O}(d)$. A specific example highlights this difference: with a weight matrix sized $128\times128$, LoRA introduces $2,048$ learnable parameters when $r'=8$, whereas OFT requires only $960$ trainable parameters at $r=8$, without any block sharing.

\subsection{Constrained Orthogonal Finetuning}

The flexibility of standard OFT can be further restricted by requiring the fine-tuned model to stay within a specified range of the pretrained weights. In COFT, the model’s forward pass is specified as

\begin{equation}\label{eq:coft_general}
    \footnotesize
    \bm{z}=\bm{W}^\top\bm{x}=(\bm{R}\cdot\bm{W}^0)^\top\bm{x},~~~\text{s.t.}~\bm{R}^\top\bm{R}=\bm{R}\bm{R}^\top=\bm{I},~\norm{\bm{R}-\bm{I}}\leq \epsilon
\end{equation}

where the finetuning transformation must be orthogonal and also close to the identity (with deviation bounded by $\epsilon$). Cayley parameterization, as described in Section~\ref{sect:eff_ortho}, enforces orthogonality. However, it is not straightforward to impose the $\epsilon$-closeness constraint using the Cayley-parameterized form. To examine the structure further, we leverage the Neumann series to approximate $\bm{R} = (\bm{I}+\bm{Q})(\bm{I}-\bm{Q})^{-1}$, yielding the first-order expansion $\bm{R} \approx \bm{I} + 2\bm{Q} + \mathcal{O}(\bm{Q}^2)$ under the assumption that the Neumann

series converges in the operator norm). Consequently, it is permissible to reformulate the constraint $\thickmuskip=2mu \medmuskip=2mu\|\bm{R}-\bm{I}\|\leq\epsilon$ within the scope of the Cayley transform, leading to an equivalent restriction of the form $\thickmuskip=2mu \medmuskip=2mu \|\bm{Q}-\bm{0}\|\leq\epsilon'$, where $\bm{0}$ represents the zero matrix and $\epsilon'$ is a distinct hyperparameter governing allowable deviation (separate from $\epsilon$). This alternative condition on $\bm{Q}$ is readily implemented via projected gradient descent. To ensure that the orthogonal matrix $\bm{R}$ is initially set to the identity, one simply initializes $\bm{Q}$ as the zero matrix. COFT can be interpreted as uniting two explicit forms of regularization: the first constrains the difference in hyperspherical energy to be minimal, while the second enforces bounded deviation from the parameters of the pretrained model. The latter form of regularization is frequently implicit in standard finetuning methodologies, but COFT introduces it as a direct, explicit component. While OFT delivers strong empirical results, we find that incorporating the explicit deviation constraint in COFT enhances finetuning stability relative to OFT. Figure~\ref{fig:eps_coft} illustrates how the value of $\epsilon$ influences COFT's performance. The plot demonstrates that $\epsilon$ mediates the trade-off between flexibility and similarity to the initial model: larger $\epsilon$ yields finetuned models more akin to those from OFT, whereas smaller values of $\epsilon$ cause the COFT-finetuned model to increasingly approximate the original pretrained text-to-image diffusion model.

\subsection{Re-scaled Orthogonal Finetuning}

We introduce a straightforward modification to standard OFT by augmenting it with a learnable scaling coefficient applied independently to each neuron. The motivation lies in the observation that such re-scaling leaves hyperspherical energy unchanged, as normalization absorbs the magnitude. More concretely, we define the forward operation as $\thickmuskip=2mu \medmuskip=2mu\bm{z}=(\bm{R}\bm{W}^0\bm{D})^\top\bm{x}$\footnote{\textbf{Errata}: The NeurIPS camera-ready paper erroneously presented the forward computation for re-scaled OFT as $\thickmuskip=2mu \medmuskip=2mu\bm{z}=(\bm{D}\bm{R}\bm{W}^0)^\top\bm{x}$. The implementation is in fact correct, ensuring the results in Appendix~\ref{app:rescaled} remain valid.} with $\thickmuskip=2mu \medmuskip=2mu \bm{D}=\text{diag}(s_1,\cdots,s_n)\in\mathbb{R}^{n\times n}$ a trainable diagonal matrix, where all diagonal entries $s_1,\cdots,s_n$ are constrained to be positive. Unlike the original OFT formulation in Eq.~\eqref{eq:oft_general}—in which $\bm{R}$ is the only trainable component—both $\bm{R}$ and $\bm{D}$ are learnable in this version. Incorporating the rescaling mechanism extends OFT’s modeling capacity, requiring only a minor increase in parameter count. Nonetheless, for the experimental results presented in the main text, we adhere to the baseline OFT formulation to better isolate the effects of orthogonal transformations; however, empirical evidence in Appendix~\ref{app:rescaled} suggests re-scaled OFT typically outperforms standard OFT.

\section{Intriguing Insights and Discussions}

\textbf{OFT accommodates arbitrary neural architectures}. In theory, OFT is applicable to all types of neural networks. While LoRA is commonly deployed on attention parameters within Transformers~\cite{hulora2022}, to maintain fair comparison with LoRA in our experiments, we likewise restrict OFT’s finetuning to the attention weights. OFT’s approach is not limited to fully connected networks; its structure makes it naturally amenable for adapting convolutional layers as well, as the block-diagonal nature of $\bm{R}$ admits unique interpretations for convolutional operations—an advantage not present for LoRA. For example, choosing as many blocks as input channels results in each block independently transforming a separate neuron channel, reminiscent of learning individual depth-wise convolution filters~\cite{chollet2017xception}. Conversely, if $\bm{R}$'s blocks are shared, the orthogonal matrix transformation is applied uniformly across all channels.

\textbf{Connection to LoRA}. LoRA introduces a low-rank modification to the weight matrix, thus safeguarding the pretrained parameters from significant perturbation. Conversely, OFT enforces that the transformation applied to the pretrained matrix is orthogonal (and thus full-rank), which ensures that the structure learned during pretraining remains largely intact. The forward computation in OFT can be equivalently expressed as $\thickmuskip=2mu \medmuskip=2mu \bm{z}=(\bm{R}\bm{W}^0)^\top\bm{x}=(\bm{W}^0+(\bm{R}-\bm{I})\bm{W}^0)^\top\bm{x}$, where $\thickmuskip=2mu \medmuskip=2mu (\bm{R}-\bm{I})\bm{W}^0$ functions similarly to the low-rank adjustment introduced in LoRA. When $\bm{W}^0$ is full-rank, OFT naturally induces a low-rank adaptation if $\thickmuskip=2mu \medmuskip=2mu \bm{R}-\bm{I}$ itself is low-rank. Furthermore, as LoRA utilizes a rank hyperparameter $r'$ to restrict the number of learned parameters, OFT analogously uses a block-diagonal size $r$ for parameter reduction. Fundamentally, OFT and LoRA embody two different paradigms of parameter-efficient learning: while LoRA relies on leveraging low-rank structures, OFT capitalizes on sparse patterns, notably block-diagonal orthogonality, to achieve compactness.

\textbf{Why OFT converges faster}. As evidenced in Figure~\ref{fig:toy_ae}, altering neuron orientation is the most potent means to effect semantic changes—precisely the type of modification that OFT facilitates. Additionally, OFT can be interpreted as fine-tuning within the hypersphere manifold, which smooths the optimization landscape and accelerates convergence. This advantage is also supported by findings in \cite{liu2021orthogonal}.

\textbf{Why not minimize hyperspherical energy}. Our approach differs from \cite{liu2021orthogonal} in that we do not target minimizing the hyperspherical energy. In classification, eliminating redundancy among neurons is preferable, but enforcing a minimum hyperspherical energy arranges neurons equidistantly over the hypersphere. This objective is inappropriate for finetuning tasks, as it risks erasing crucial information encoded during pretraining.

\textbf{Balancing flexibility and regularity in finetuning}. Our analysis reveals an inherent tension between achieving flexibility and maintaining regularity during model finetuning. Traditional finetuning approaches offer maximal flexibility, yet often suffer from instability and are prone to model degradation. In contrast, the straightforward OFT method manages to effectively reconcile this trade-off by maintaining pairwise neuron angles, thus striking a favorable equilibrium. Additionally, the block-diagonal formulation can be interpreted as imposing stricter regularization on the orthogonal matrix.

\textbf{Zero inference-time overhead}. In contrast to frameworks like ControlNet, the OFT method incurs no extra computational cost during inference. Once training is complete, inference only requires multiplying the learned orthogonal matrix $\bm{R}$ by the original weight matrix $\bm{W}^0$, producing a functionally equivalent matrix $\thickmuskip=2mu \medmuskip=2mu \bm{W}=\bm{R}\bm{W}^0$. This process preserves inference speed identical to that of the pretrained model.

\section{Experiments and Results}

\textbf{Experimental protocol}. Our experiments utilize Stable Diffusion v1.5~\cite{rombach2022high} as the foundation text-to-image generator. To avoid selection bias, output images for each technique are sampled randomly. For evaluations in subject-driven generation, procedures largely mirror those in DreamBooth~\cite{ruiz2023dreambooth}; for controlled generation, protocols are adapted from ControlNet~\cite{zhang2023adding} and T2I-Adapter~\cite{mou2023t2i}. To facilitate fair comparisons against LoRA, OFT or COFT is only applied to the same network layer targeted by LoRA. Supplementary information, including additional findings and methodological details, is presented in Appendix~\ref{app:settings}.

\subsection{Subject-driven Generation}

\textbf{Configuration}. We compare against DreamBooth~\cite{ruiz2023dreambooth} and LoRA~\cite{hulora2022} as baseline approaches. All models use the same loss as described in the DreamBooth framework. The training and hyperparameter setups for DreamBooth and LoRA generally adhere to their respective publications, using their recommended best configurations. Expanded results can be found in Appendices~\ref{app:settings},\ref{app:coft_oft},\ref{app:more_qualitative}, and~\ref{app:failure}.

\textbf{Finetuning stability and convergence}. Our initial analysis focuses on the stability and convergence rate of finetuning across DreamBooth, LoRA, OFT, and COFT, as illustrated in Figure~\ref{fig:teaser} and Figure~\ref{fig:db_convergence}. The findings indicate that both OFT and COFT demonstrate robust finetuning dynamics for the diffusion model, maintaining stability throughout the process. At the 400-iteration mark, DreamBooth and OFT-based methods succeed in controlling image generation effectively, whereas LoRA does not manage to retain subject identity. Progressing to 2000 iterations, DreamBooth encounters mode collapse—producing degenerate images—and LoRA further deteriorates, failing to represent key subject features (e.g., confusing yellow shirt with yellow fur). By contrast, OFT and COFT continue to exhibit stable control and consistency in the output images. This evidence substantiates the efficacy of the OFT method in achieving rapid and steady convergence for subject-specific image synthesis. Moreover, the low sensitivity of OFT and COFT to the number of finetuning iterations simplifies practical deployment, as it reduces the need to tune this hyperparameter individually for different subjects. Both approaches permit the direct specification of a large iteration count without the necessity for extensive parameter searching. Notably, with COFT and an appropriate choice of $\epsilon$, the process of selecting learning rates and iteration numbers becomes trivial.

\textbf{Quantitative comparison}. Consistent with the evaluation setup in \cite{ruiz2023dreambooth}, we present a quantitative assessment encompassing subject fidelity (measured by DINO~\cite{caron2021emerging} and CLIP-I~\cite{radford2021learning}), adherence to the text prompt (CLIP-T~\cite{radford2021learning}), and generation diversity (LPIPS~\cite{zhang2018unreasonable}). The CLIP-I metric calculates the mean pairwise cosine similarity between CLIP embeddings derived from generated images and those from real images of the subject. Similarly, DINO applies the same procedure using ViT S/16 DINO embeddings. The CLIP-T metric reflects the mean cosine similarity between embeddings of the text prompts and the corresponding generated images. Generation diversity is gauged through the average LPIPS cosine similarity across samples generated for a given subject and prompt. As reported in Table~\ref{table:dreambooth_final}, COFT and OFT both achieve significant improvements over DreamBooth and LoRA on the DINO and CLIP-I scores, while maintaining equivalent or slightly improved performance on prompt fidelity and diversity criteria. To reduce variability, we repeat the finetuning process 30 times for each method using different random seeds on the same pretrained model. Overall, these experiments demonstrate that the OFT framework not only converges faster and more reliably but also consistently delivers superior final outcomes.

\textbf{Qualitative comparison}. To facilitate a clearer appreciation of OFT's advantages, we present several randomly selected cases of subject-driven generation in Figure~\ref{fig:dreambooth_final}. To maintain experimental fairness, each method is applied to generate outputs from the identical finetuned model, ensuring that no text prompt is specially adapted to favor a particular result. For every approach, we choose the model version that records the highest validation CLIP score. The visual results displayed in Figure~\ref{fig:dreambooth_final} suggest that both OFT and COFT are able to robustly retain semantic consistency with the intended subject, while LoRA frequently fails to do so (for instance, LoRA does not preserve the subject’s characteristics in the bowl scenario). Furthermore, OFT and COFT exhibit a much greater capacity for prompt-based control, whereas DreamBooth, although it generally maintains subject identity, regularly does not align the generation with the textual instruction (such as in the top row of the bowl example). Collectively, this qualitative analysis indicates that our OFT framework excels in balancing subject fidelity with prompt controllability. In addition, the performance of OFT is relatively unaffected by the number of iterations used; OFT maintains strong outcomes even when iteration counts are high, a property not shared by either DreamBooth or LoRA. Additional qualitative samples are available in Appendix~\ref{app:more_qualitative}. Also, a human evaluation study, as detailed in Appendix~\ref{app:human}, offers further confirmation of OFT’s effectiveness.

\subsection{Controllable Generation}

\textbf{Settings}. For our baselines, we consider ControlNet~\cite{zhang2023adding}, T2I-Adapter~\cite{mou2023t2i}, and LoRA~\cite{hulora2022}. The primary manuscript investigates three notably challenging controllable generation tasks: canny edge to image~(C2I) on the COCO dataset~\cite{lin2014microsoft}, segmentation map to image~(S2I) on the ADE20K dataset~\cite{zhou2017scene}, and landmark to face~(L2F) on the CelebA-HQ dataset~\cite{karras2018progressive,xia2021tedigan}. Each method is used to finetune Stable Diffusion~(SD) v1.5 for 20 epochs across these three datasets. Supplementary results and further task-specific evaluations can be found in Appendix~\ref{app:more_qualitative},~\ref{app:more_control_task},~\ref{app:failure}.

\textbf{Convergence}. We investigate the convergence characteristics of ControlNet, T2I-Adapter, LoRA, and COFT within the L2F benchmark. Our analysis includes both quantitative and qualitative assessments. For our primary metric, we utilize the average $\ell_2$ norm between the predicted and the target facial landmarks. As depicted in Figure~\ref{fig:face_convergence}, the facial landmark error is tracked across different numbers of fine-tuning epochs for each model. It is evident from the results that COFT and OFT display markedly quicker convergence rates. Notably, LoRA requires 20 epochs to attain a performance level comparable to OFT at merely the 8-th epoch. Both OFT and COFT operate with a similar number of trainable parameters as LoRA—significantly fewer than ControlNet—but outperform existing baselines in terms of convergence efficiency. Moreover, OFT's rapid convergence is corroborated by the qualitative examples shown in Figure~\ref{fig:teaser}, where the rightmost illustration demonstrates OFT's higher data efficiency; it converges robustly with only 5\% of the ADE20K training set, outperforming both ControlNet and LoRA under the same conditions. For qualitative assessment, we center our comparison on OFT, COFT, and ControlNet, since ControlNet achieves the nearest landmark error to our approaches. Figure~\ref{fig:face_convergence_qual} illustrates that OFT and COFT not only converge steadily but also progressively align generated facial poses to the specified control landmarks. In contrast, ControlNet’s control capability exhibits an abrupt onset post 8-th epoch, mirroring the sudden convergence behavior noted in \cite{zhang2023adding}. This enhanced stability of convergence in our OFT framework renders it more practical, with smoother and more foreseeable training dynamics, thereby facilitating hyperparameter selection.

\textbf{Quantitative comparison}. To assess the effectiveness of controllable generation, we propose a metric for control consistency. This metric involves deriving the control signal from the output image and comparing it to the original input control signal. Specifically, for the C2I task, we report IoU and F1 score; for S2I, we report mean IoU along with both mean and overall accuracy; and for L2F, we use the average $\ell_2$ distance between the true and predicted control landmarks. Appendix~\ref{app:settings} contains further details regarding these consistency metrics. Each method evaluated was tuned with optimal hyperparameter selections. According to Table~\ref{table:control_gen}, both OFT and COFT demonstrate superior and more precise control in comparison with alternative methods. Notably, adapter-based strategies such as T2I-Adapter and ControlNet exhibit slower convergence and ultimately lower performance. LoRA achieves higher results than ControlNet in the S2I scenario, but falls behind on C2I and L2F tasks. Generally, even after careful hyperparameter optimization, LoRA displays a lower upper-bound on performance. Conversely, our observations indicate that the OFT framework continues to improve with increasing trainable parameters, suggesting its full capacity has not yet been reached. Importantly, the quantitative assessment protocol we adopt for controllable generation represents a novel aspect of our work, as it provides an accurate measure of control performance for finetuned models—subject to the precision limits of the external segmentation or detection models employed.

\textbf{Qualitative comparison}. We further present a qualitative evaluation of OFT and COFT against leading approaches such as ControlNet, T2I-Adapter, and LoRA. As illustrated in the random samples of Figure~\ref{fig:control_final}, both OFT and COFT deliver images with high visual fidelity and realism, alongside precise control over the outputs. For the S2I task, it is evident that LoRA is unable to generate images conforming to the provided segmentation maps, whereas ControlNet, OFT, and COFT all manage to effectively guide image synthesis according to the input. Distinct from ControlNet, both OFT and COFT generate images that are richer in detail and more plausible in terms of geometric arrangement, yet require significantly fewer parameters. In the C2I task, OFT and COFT both succeed in producing semantically meaningful results from coarse Canny edge maps, in contrast to the poorer performance observed with T2I-Adapter and LoRA. Regarding the L2F task, our approach yields the most precise pose-controlled face generations, maintaining accuracy even with complex facial orientations. Across all three evaluated tasks, the qualitative outputs from OFT and COFT surpass those of existing state-of-the-art methods, highlighting the advantages of our OFT framework for controllable image generation. For additional qualitative results, we supply further examples covering all three tasks in Appendix~\ref{app:control_exp}. Furthermore, in Appendix~\ref{app:more_control_task}, we show that OFT generalizes effectively to a broader set of control challenges, including dense pose to human body, sketch to image, and depth to image tasks.

\section{Concluding Remarks and Open Problems}

Building on the insight that angular relationships among neurons play a pivotal role in shaping visual semantics, we introduce a straightforward yet robust finetuning technique—orthogonal finetuning—for manipulating text-to-image diffusion models. Our approach focuses on two core tasks: subject-driven generation and controllable generation. OFT achieves superior controllability and maintains greater finetuning stability, all while requiring a reduced set of tunable parameters in comparison to established methods. Crucially, OFT incurs no extra computational cost during inference, making it a practical choice for deployment.

Nevertheless, OFT gives rise to several compelling open questions. The approach enforces orthogonality using Cayley parametrization, which necessitates matrix inversion—an operation that somewhat restricts the scalability of OFT. While our block diagonal parametrization partially mitigates this constraint, developing faster, differentiable solutions for matrix inversion is still unresolved. Moreover, OFT is distinct in its potential for compositionality: orthogonal matrices from multiple OFT finetuning processes can be multiplied and the resulting matrix remains orthogonal. Investigating whether these combined matrices collectively encode the expertise of all respective downstream tasks is an intriguing avenue for further research. Lastly, OFT's parameter efficiency is closely tied to the block diagonal structure, which inevitably brings about certain biases and limits adaptability. Identifying strategies to enhance parameter efficiency without introducing such biases stands as a significant open challenge.